\cleardoublepage
\vspace*{2\baselineskip}
{\HduLevelOneTitleFormat 摘要\par}
\vspace{2\baselineskip}

{\song\zihao{-4}\setlength{\baselineskip}{20pt}
大语言模型在代码生成，与错误诊断上的可用性，已获验证。自然语言交互、个性化反馈，两端尤为突出。当前编程教学，在资源可得性、辅导时效、过程追踪，和即时反馈上，仍有短板。入门阶段，兴趣滑坡并不少见。本文据此，设计了一套面向中小学生的智能编程学习平台，在因材施教与教学效率之间，寻找折中。

架构上前后端分离。交互界面，基于 Nuxt 3 与 Vue 3 搭建，样式依托 Tailwind CSS、DaisyUI。核心业务，由 FastAPI、SQLAlchemy 承担。数据层选用 SQL Server 持久化，Redis 缓存加速，阿里云 OSS 存放教学资源。

功能覆盖用户认证、课程编排、在线编程实践、任务管理、师生互动，及教学资源管理。大模型驱动的辅学模块，承接自由问答、代码逻辑拆解、编译错误分析、问题分解引导，以及抽象建模与算法设计指导。针对中小学生认知节奏，提示词策略、输出语气，和错误反馈均做了调整，反馈逐步给出，而非一次倾倒。

全文围绕需求分析、总体设计、核心模块实现与功能验证展开。平台已为学生日常编程练习，与教师教学管理，提供较为完整的操作支撑，在代码稳定运行、错误反馈、进度跟踪与资源组织方面，已有初步实践基础。后续可在 RAG 私域知识增强、认知模板评估，及更细粒度的个性化推荐方向，继续拓展。

\vspace{20pt}
\noindent{\CJKfontspec[FakeBold=2.5]{Heiti SC}\zihao{-4}关键词：}\song\zihao{-4}大语言模型；编程教育；智能辅学；在线编程；教学平台
}

\cleardoublepage
\vspace*{2\baselineskip}
{\HduLevelOneTitleFormat ABSTRACT\par}
\vspace{2\baselineskip}

{\fontspec{Times New Roman}\zihao{-4}\setlength{\baselineskip}{20pt}
The availability of large language models in code generation and error diagnosis has been verified. Natural language interaction and personalized feedback are particularly prominent at both ends. In the current programming teaching, there are still shortcomings in resource availability, counseling timeliness, process tracking, and instant feedback. In the introductory stage, interest landslides are not uncommon. Based on this, this paper designs a set of intelligent programming learning platform for primary and middle school students, and finds a compromise between teaching students in accordance with their aptitude and teaching efficiency. 
 
The front and back ends are separated on the architecture. The interactive interface is built based on Nuxt 3 and Vue 3, and the style relies on Tailwind CSS and DaisyUI. The core business is undertaken by FastAPI and SQLAlchemy. The data layer uses SQL Server persistence, Redis cache acceleration, and Alibaba Cloud OSS to store teaching resources. 
 
The functions cover user authentication, course arrangement, online programming practice, task management, teacher-student interaction, and teaching resource management. The large model-driven auxiliary module undertakes free question answering, code logic disassembly, compilation error analysis, problem decomposition guidance, and abstract modeling and algorithm design guidance. In view of the cognitive rhythm of primary and middle school students, the prompt word strategy, output tone, and error feedback have been adjusted, and the feedback is gradually given, rather than a dumping. 
 
The full text focuses on demand analysis, overall design, core module implementation and functional verification. The platform has provided a relatively complete operational support for students ' daily programming exercises and teachers ' teaching management. It has a preliminary practical basis in terms of stable code operation, error feedback, progress tracking and resource organization. Subsequently, it can continue to expand in the RAG private domain knowledge enhancement, cognitive template evaluation, and more fine-grained personalized recommendation direction.

\vspace{20pt}
\noindent\fontspec{Times New Roman}\zihao{-4}\textbf{Key Words：}\fontspec{Times New Roman}\zihao{-4}large language model; programming education; intelligent tutoring; online programming; learning platform
}
